% docs/manual/operations/chapters/06-stations-at-scale.tex
% 第 6 章：大规模站点评估实战

\chapter{大规模站点评估实战}

\openbench{} 的站点评估有独特挑战：每站独立、IO 密集、对 HDF5 lock 敏感。本章给"千站点级"评估的运维要点。

\section{站点评估的特殊性}

与 grid 评估对比：

\begin{longtable}{p{3cm} p{4.5cm} p{5cm}}
  \toprule
  指标 & Grid & Station \\
  \midrule
  \endhead
  并行度 & 单节点 & 站点级（实际并行） \\
  IO 模式 & 大文件少次访问 & 大量小文件多次访问 \\
  内存 & 与分辨率相关 & 与站点数相关 \\
  HDF5 lock 风险 & 低 & 高 \\
  典型时长 & 几十分钟 & 几小时 \\
  \bottomrule
\end{longtable}

\section{Fulllist CSV 格式}

站点评估必备：站点列表 CSV。

\begin{minted}{text}
site,lat,lon,IGBP
US-Ha1,42.5378,-72.1715,DBF
US-MMS,39.3232,-86.4131,DBF
DE-Tha,50.9624,13.5651,ENF
...
\end{minted}

\textbf{必填列}：\texttt{site}、\texttt{lat}、\texttt{lon}。其他列被忽略但常含 IGBP / 海拔等元信息便于过滤。

\subsection{多源 fulllist 合并}

如果你要同时评估 FLUXNET + AmeriFlux + EuroFlux：

\begin{minted}{python}
import pandas as pd
dfs = [pd.read_csv(p) for p in ["fluxnet.csv", "amf.csv", "eflux.csv"]]
all_sites = pd.concat(dfs).drop_duplicates(subset="site")
all_sites.to_csv("merged.csv", index=False)
\end{minted}

\section{station\_scanner：从 NC 批量扫站点}

\modname{data.station\_scanner} 从 station NC 文件中提取所有站点：

\begin{minted}{python}
from openbench.data.station_scanner import scan_stations

sites = scan_stations("/data/PLUMBER2/")
# DataFrame: site, lat, lon, available_variables, time_range
\end{minted}

输出可作为 fulllist 输入；典型用于"我有 NC 数据但还没列表"。

\section{station\_matcher：grid sim → 站点采样}

当 simulation 是 grid 模型时，\modname{data.station\_matcher} 在 grid 中按站点位置采样：

\begin{itemize}
  \item 默认：最近邻（nearest grid cell）
  \item 备选：双线性（4 邻接点加权）
  \item 落空：站点位置在 grid range 之外（lat/lon 不覆盖） → 跳过该站
\end{itemize}

\subsection{落空站点诊断}

\begin{minted}{bash}
# 看 run.log 中的 skipped sites
grep -i "skipping station" output/<run>/run.log

# 典型原因
# - 站点 lat/lon 在 lat_range 之外
# - 站点位置全为 NaN（海洋 / 极地）
# - sim NC 该位置 mask 掉了
\end{minted}

\section{千站点资源预算}

\subsection{时间预算}

理想情形：单站约 5 秒（小变量、月尺度）。

\begin{itemize}
  \item 100 站 → 500 秒（8 分钟）单核
  \item 1000 站 → 80 分钟单核 / 5 分钟 16 核
  \item 10000 站 → 13 小时单核 / 50 分钟 16 核
\end{itemize}

\subsection{内存预算}

每站 ~10-50 MB（含 sim 切片 + ref 数据 + 中间）。1000 站 + 16 worker = 800 MB / worker = 13 GB 总。

\subsection{磁盘}

每站对应一个 metric/score CSV 行；千站量级是 KB-MB 级，不是瓶颈。

\subsection{IO 频次}

每站要读 sim NC 一个切片 + ref 一个时间序列。1000 站 × 16 worker = 16000 次 NC 读。这是 HDF5 lock 风险最大的地方。

\section{HDF5 locking：历史与现状}

2026-04-03 之前 \openbench{} 在大规模站点评估上经常崩溃：

\begin{minted}{text}
OSError: [Errno 11] Resource temporarily unavailable
HDF5: file lock
\end{minted}

\subsection{根因}

HDF5 1.14+ 默认对 NC 文件加 advisory lock，多 worker 同时只读访问触发。

\subsection{修复}

主流程在 \cli{run} 启动时设：

\begin{minted}{python}
os.environ.setdefault("HDF5_USE_FILE_LOCKING", "FALSE")
\end{minted}

\begin{warnBox}
\textbf{Notebook 自定义脚本要点}：必须在 \texttt{import openbench} 或 \texttt{import xarray}（任一引入 HDF5 库的语句）\emph{之前} 设这个变量，否则 HDF5 已经把 \texttt{TRUE} 缓存进 process。

\begin{minted}{python}
import os
os.environ["HDF5_USE_FILE_LOCKING"] = "FALSE"
# 现在再 import
from openbench.runner.local import run_evaluation
\end{minted}
\end{warnBox}

\subsection{2026-04-03 修复全集}

memory 中记录的"Station evaluation fixes"还包括：

\begin{itemize}
  \item nested parallelism 死锁（joblib × dask 嵌套）
  \item ref symlink 处理
  \item fallback variable 的单位换算
  \item None tolerance（某些字段为 None 时的兜底）
\end{itemize}

具体修复见 git log 与 \file{docs/superpowers/reviews/} 历史 audit。

\section{失败站点的批量诊断}

\subsection{看哪些站点失败}

\begin{minted}{bash}
grep -A 1 "ERROR.*station" output/<run>/run.log | head -50
# 或
grep "skipping station" output/<run>/run.log | wc -l
\end{minted}

\subsection{失败模式分类}

\begin{longtable}{p{4cm} p{4cm} p{5cm}}
  \toprule
  失败模式 & 提示 & 应对 \\
  \midrule
  \endhead
  位置出范围 & "out of lat\_range" & 扩大 \yamlkey{lat\_range} \\
  Sim NaN & "sim all NaN at site" & sim 该位置无数据 \\
  Ref 时段不够 & "insufficient overlap" & 调 \yamlkey{years} 或 \yamlkey{min\_year\_threshold} \\
  HDF5 lock & "Resource temporarily" & 检查 env var \\
  Station NC 损坏 & "Cannot open file" & 重新下载该站数据 \\
  \bottomrule
\end{longtable}

\section{实战配置示例：FLUXNET 全球}

\begin{minted}{yaml}
project:
  name: fluxnet-global-eval
  output_dir: /scratch/$USER/fluxnet-eval
  years: [2003, 2014]
  tim_res: Hour
  num_cores: 32
  time_alignment: per_pair
  min_year_threshold: 3

evaluation:
  variables:
    - Latent_Heat
    - Sensible_Heat
    - Net_Radiation
    - Gross_Primary_Productivity

reference:
  Latent_Heat: FLUXNET_PLUMBER2
  Sensible_Heat: FLUXNET_PLUMBER2
  Net_Radiation: FLUXNET_PLUMBER2
  Gross_Primary_Productivity: FLUXNET_PLUMBER2

simulation:
  CoLM2024:
    model: CoLM2024
    root_dir: /shared/sim/CoLM2024
    fulllist: /shared/PLUMBER2/list/PLUMBER2_sites.csv
\end{minted}

\textbf{运行}：

\begin{minted}{bash}
sbatch <<EOF
#!/bin/bash
#SBATCH --cpus-per-task=32
#SBATCH --mem=64G
#SBATCH --time=12:00:00

source ~/miniforge3/bin/activate openbench
openbench run fluxnet-global.yaml --cores \$SLURM_CPUS_PER_TASK
EOF
\end{minted}

预期：~200 站 × 4 变量 = 800 个 (var × ref) 评估单元，~3-5 小时跑完。

\section{下一步}

下一章详解监控与日志：logging\_system 详解、外部监控集成、长任务监控。
